\begin{salyga}
Su kuriomis parametro $k$ reikšmėmis lygtis turi lygiai vieną realųjį sprendinį?
    \[ (k - 1)x^2 - 2(k + 1)x + k + 4 = 0 \]
\end{salyga}

\begin{atsakymas}
$k=1$ arba $k=5$.
\end{atsakymas}

\begin{sprendimas}
Jei $k=1$, lygtis tampa tiesine $-4x+5=0$, todėl turi vieną sprendinį $x=\frac54$.

Jei $k\ne1$, lygtis kvadratinė ir turi lygiai vieną realųjį sprendinį tada ir tik tada, kai
\[
\frac D4=(k+1)^2-(k-1)(k+4)=5-k=0.
\]
Gauname $k=5$, ir ši reikšmė tenkina sąlygą $k\ne1$. Vadinasi, tinka $k=1$ ir $k=5$.
\end{sprendimas}
